\begin{result}[Implication of CA]
    $x^2 = 2, x > 0$ has a unique solution in $\mathbb{R}$, called $\sqrt{2}$
\end{result}

\begin{proof}
    Let $S:= \{ x \in \mathbb{R}, x > 0, x^2 < 2\} $. 
    We claim that 
    \begin{itemize}
        \item $S \neq \emptyset$. Take $1, 1 > 0, 1^2 < 2$. Check.
        \item $S$ is bounded from above. Take $2$. Check.
Suppose $\exists x \in S, x > 0$ then $x^2 > 2$. 
    \end{itemize} 
    Hence by CA, $S$ has a supremum in $\mathbb{R}$, 
    say $m$. 
    \medskip 

    Claim: $m > 0, m^2 = 2$. It suffices to check $m^2 = 2$.
    ATC otherwise.
    \begin{enumerate}
        \item If $m^2 < 2$, 
then we can find $n \in \mathbb{N}$
s.t. $ \left( m + \frac{1}{n} \right)^2 < 2 $.
\medskip

By AP, $\exists n \in \mathbb{N}$ s.t. 
$n > \frac{2m + 1}{2-m^2}$, then 
\begin{align*}
    n > \frac{2m + 1}{2 - m^2} & \implies \frac{2m + 1}{n} < 2 - m^2 \\
    & \implies m^2 + \frac{2m}{n} + \frac{1}{n} < 2 - m^2 \\
    & \implies m^2 + \frac{2m}{n} + \frac{1}{n} < 2 \\
    & \implies m^2 + \frac{2m}{n} + \frac{1}{n^2} < m^2 + \frac{2m}{n} + \frac{1}{n} < 2 \\
    & \implies \left( m + \frac{1}{n} \right)^2 < 2
\end{align*}
        \item Assume $m^2 > 2$, then we can find $n \in \mathbb{N}$ 
such that $m - \frac{1}{n}$ is also an upper bound of $S$.
\medskip 

It suffices to find an $n$ such that $ \left( m - \frac{1}{n} \right)^2 > 2 $, 
which implies $ \left( m - \frac{1}{n} \right)^2 > x^2 \implies m - \frac{1}{n}> x$, 
$\forall x \in S$. 
\medskip 

By AP, $\exists n \in \mathbb{N}$ s.t. 
$n > \frac{2m}{m^2 - 2}$, then 
\begin{align*}
    n > \frac{2m}{m^2 - 2} & \implies m^2 - 2 > \frac{2m}{n} \\
    & \implies m^2 - \frac{2m}{n} > 2 \\
    & \implies m^2 - \frac{2m}{n} + \frac{1}{n^2} > m^2 - \frac{2m}{n} > 2 \\
    & \implies  \left(  m - \frac{1}{n} \right)^2 > 2
\end{align*}
    \end{enumerate} 
\end{proof}

\begin{definition}[Density]
    $A \subset B$ is dense if $
    \forall x \in B, \forall \epsilon > 0,
    \exists y \in A 
    $ s.t. $ |x - y| < \epsilon$.
    \medskip
\end{definition}
\begin{remark}
    We say $A$ is dense in $B$ if 
    you can do $\epsilon$-approximation for anything in $B$ 
    with something in $A$.
\end{remark}



\begin{theorem}
    We have
    \begin{enumerate}
        \item $\mathbb{Q}$ is dense in $\mathbb{R}$.
        \item $\mathbb{R} \setminus \mathbb{Q}$ is dense in $\mathbb{R}$.
    \end{enumerate}
\end{theorem}

\begin{remark}
    $\mathbb{Q}$ dense in $\mathbb{R}$ 
    $\iff \forall a \in \mathbb{R}, b \in \mathbb{R}$, 
    $a < b$, 
    $\exists r \in \mathbb{Q}$ s.t. 
    $a < r < b$.
\end{remark}

\begin{remark}
    $r \in \mathbb{Q}, r \neq 0 \implies \sqrt{2}r \in \mathbb{R} \setminus \mathbb{Q}$
\end{remark}

\begin{proof}[Proof. Thm 1.19]
    \begin{enumerate}
        \item 
    We first show that $(1) \implies (2)$, i.e. $\mathbb{Q}$ dense $\implies$ 
    ${R} \setminus \mathbb{Q}$ dense. 
    \medskip 

    $\forall x \in \mathbb{R}, \forall \epsilon > 0$. WTS 
    $\exists y \in \mathbb{R} \setminus \mathbb{Q}$ s.t. $|x - y| < \epsilon$. 
    \medskip

    Let $x^{\prime} = \frac{x}{\sqrt{2}}$, 
    $\epsilon^{\prime} = \frac{\epsilon}{\sqrt{2}}$. 
    Since $\mathbb{Q}$ dense in $\mathbb{R}$, 
    $\exists r \in \mathbb{Q}$ s.t. $r \neq 0, \left| x^{\prime} - r \right| < \epsilon^{\prime}$. (If $r = 0$, then we can pick a smaller window $\epsilon^{\prime \prime} = \frac{r - x^{\prime}}{2}$.)

    \begin{center}
    \begin{tikzpicture}[x=1.4cm]
    \draw[-{Stealth[length=2.2mm]}] (-0.4,0) -- (5.4,0) node[right] {$\mathbb{R}$};
    \foreach \p/\lbl in {1.5/{x'-\ep'}, 3/{x'}, 2.4/{r}, 4.5/{x'+\ep'}}
        \filldraw (\p,0) circle (1.6pt) node[above=4pt] {$\lbl$};
    \end{tikzpicture}
    \end{center}


    Then, we take $y = \sqrt{2} r$. 
    \begin{align*}
        |x - \sqrt{2} r| &= \sqrt{2} |x^{\prime} - r| \\
        & < \sqrt{2} \epsilon^{\prime} \\
        & = \epsilon
    \end{align*}

    \item 
Now we show (1), i.e. $\mathbb{Q}$ is dense in $\mathbb{R}$. 
I.e. WTS $\forall x \in \mathbb{R}, \forall \ep > 0$
$\exists y \in \mathbb{Q}$ s.t. 
$|x - y | < \ep$. 
\medskip 

By Archimedean property, $\exists n \in \mathbb{N} \text{s.t.} n \ep > 1$
\medskip

\begin{center}
\begin{tikzpicture}[x=1.3cm]
  % original scale
  \draw[-{Stealth[length=2mm]}] (-2.3,1.3) -- (2.3,1.3) node[right] {$\mathbb{R}$};
  \foreach \p/\lbl in {-1/{x-\ep}, 0/{x}, 1/{x+\ep}}
    \filldraw (\p,1.3) circle (1.5pt) node[above=3pt] {$\lbl$};

  \node[font=\small] at (2.9,0.65) {$\times\, n$};

  % blown up by n
  \draw[-{Stealth[length=2mm]}] (-2.9,0) -- (2.9,0) node[right] {$\mathbb{R}$};
  \foreach \p/\lbl in {-2.4/{nx-n\ep}, 0/{nx}, 2.4/{nx+n\ep}}
    \filldraw (\p,0) circle (1.5pt) node[above=3pt] {$\lbl$};
  \filldraw[red] (-0.9,0) circle (1.5pt) node[above=3pt] {$m$};
  \draw[{Stealth[length=1.6mm]}-{Stealth[length=1.6mm]}] (0,-0.65) --
    node[fill=white,inner sep=1pt,font=\small] {$n\ep>1$} (2.4,-0.65);
\end{tikzpicture}
\end{center}



Also $
\forall y \in \mathbb{R}, \exists m \in \mathbb{Z}
\text{s.t.} |y - m| < 1
$ 
Hence take $y = nx$
\begin{align*}
& |nx - m | < 1 < n \ep \\
\implies & \left| x - \frac{m}{n} \right|  < \frac{n \ep}{n} = \ep
\end{align*}
    \end{enumerate}
\end{proof}

\newpage

\phantomsection
\section*{Side note: Cardinality}
\addcontentsline{toc}{section}{Side note: Cardinality}

\begin{definition}[Finite]
    A set $S \neq \emptyset$ is \textbf{finite}
    if $\exists n \in \mathbb{N}$
    and a bijection 
    $f: S \to \{ 1, 2, \hdots n \} $. 
    Otherwise $S$ is \textbf{infinite}.
\end{definition}

\begin{definition}[Countability]
    \hfill
    \begin{enumerate}
        \item A set $S$ is 
\textbf{countably infinite}
if there is a bijection $f: S \to \mathbb{N}$
        \item A set $S$ is \textbf{countable} 
if it is empty
or countably infinite 
        \item A set $S$ is \textbf{uncountable} if
it is not countable.
    \end{enumerate}
\end{definition}

\begin{remark}
    Every nonempty finite subset of $\mathbb{R}$ has a max and min. 
\end{remark}

\begin{proposition}
    \hfill 
    \begin{enumerate}
        \item Any subset of a countable set is countable.
        \item A countable union of countable sets is countable.
        \item $\mathbb{Q}$ is countable. \textbf{Exercise}
        \item $\mathbb{R} \setminus \mathbb{Q}$ is uncountable $\iff$ $\mathbb{R}$ is uncountable. \textbf{Exercise}
        \item Union of a countable set and an uncountable set is uncountable. 
    \end{enumerate}
\end{proposition}

\begin{proof}[Proof that $\mathbb{Q}$ is countable]
    There is an injection from $\mathbb{Q}$  to the set $A \coloneq \{ (a, b): a, b \in \mathbb{Z} \} $. 
    Fix arbitrary $a \in \mathbb{Z}$, $A_a \coloneq \{ (a, b) : b \in \mathbb{Z} \} $. There is a bijection 
    between $A_a$ and $\mathbb{Z}$, and $\mathbb{Z}$ countable, hence $A_a$ is countable. 
    $A = \bigcup_{a \in \mathbb{Z}}A_a$ is a countable union of countable sets and is therefore countable.
\end{proof}

\newpage 

\section{Metric spaces and point-set topology}

\subsection{Metric spaces, open and closed sets}

\begin{definition}[Metric space]
    A \textbf{metric space} $(X, d)$ is a set $X$
    of elements (called points)
    equipped with a metric or function
    $d: X \times X \to [0, \infty)$ 
    such that 
    \begin{enumerate}
        \item $\forall x, y \in X, d(x,y) \geq 0$. 
i.e. The image of $d$ is in $\mathbb{R}$ and non-negative 
        \item $d(x,y) = 0 \iff x = y$
        \item $\forall x, y \in X, d(x,y) = d(y,x)$
        \item $\forall x, y, z, \in X$, 
$d(x,z) \leq d(x, y) + d(y, z)$
    \end{enumerate}
\end{definition}

\begin{definition}[Discrete metric]
    The \textbf{discrete metric} $d$ is defined such that $\forall x, y$
    \[
        d(x,y) = \mathbf{1}[x \neq y] = 
        \begin{cases}
            0, \text{ if } x = y \\
            1, \text{ if } x \neq y 
        \end{cases}
    \]
\end{definition}


\begin{example}
    The following are metric spaces 
    \begin{itemize}
        \item $X = \mathbb{R}, d(x,y) = |x-y|$
        \item 
$
    X = \mathbb{R}^n := \{ 
        (x_1, x_2, \hdots x_n) : x_i \in \mathbb{R}
    \} 
$, 
$
d_1(x, y) = \left( \sum\limits_{i = 1}^{n} |x_i - y_i|^2 \right)^{ \frac{1}{2}}
$, i.e. Euclidean distance
        \item 
$
    X = \mathbb{R}^n 
$, 
$
d_2(x, y) = \max_i |x_i - y_i|
$
\item $X$ is the set of continuous functions on $[a, b]$, and $d(f,g) = \sup_{x \in [a, b]} |f(x) - g(x)|$
\item Any $X \neq \emptyset$ with the discrete metric.
    \end{itemize} 
\end{example}

\begin{definition}[Neighborhood]
    Let $(X,d)$ be a metric space. 
    Let $x \in X, \ep> 0$. 
    The \textbf{$\ep$-neighborhood} of $x$ 
    is 
    \[
        N(x, \ep) = B_{\ep}(x) := 
        \{ 
            y \in X, d(x, y) < \ep
        \} 
    \]
    A \textbf{deleted neighborhood} is 
    \begin{align*}
        N^{*}(x, \ep) :&= \{ 
            y \in X, 0 < d(x, y) < \ep
        \} \\
        &= N(x, \ep) \setminus \{ x \} 
    \end{align*}
\end{definition}

\begin{definition}[Interior]
    Let $(X, d)$ be a metric space, 
    $S \subset X$. 
    A point $x \in S$ is an 
    \textbf{interior point}  of $S$ if 
    \[
    \exists \ep > 0 
    \text{ s.t. }
    N(x, \ep) \subset S 
    \]

    I.e. If an $\ep$-neighborhood around $x$ 
    is in $S$, then $x$ is an interior point.

    \medskip
    
    The set of all interior points of $S$
    is called the \textbf{interior} of $S$, 
    denoted $\operatorname{int} S$
    \[
        \operatorname{int}  S := \{ 
        x \in S: \exists \ep > 0 \text{ s.t. }
        N(x, \ep) \subset S
        \} 
    \]
\end{definition}

\begin{definition}[Boundary]
    Let $(X, d)$ be a metric space, 
    $S \subset X$. 
    A point $x \in S$ is a \textbf{boundary point} of $S$ if 
    $\forall \ep > 0$,
    \begin{align*}
        N (x, \ep) \cap S \neq \emptyset \\
        N (x, \ep) \cap S^c \neq \emptyset
    \end{align*}
    The set of boundary pints of $S$ is $\operatorname{bd} S$
    \begin{align*}
        \operatorname{bd} S :&= \{ 
            x \in X: 
            \forall \ep > 0, 
            N(x, \ep) \cap S \neq \emptyset,
            N(x, \ep) \cap S^c \neq \emptyset
        \}  \\
        &= \left( \operatorname{int} S \cup \operatorname{int} S^c \right)^c
    \end{align*}
\end{definition}

\begin{example}
    $\operatorname{int} \mathbb{N} = \emptyset$, 
\end{example}
\begin{proof}
    Fix $x \in \mathbb{R}$ and $\varepsilon > 0$ arbitrary. 
    Take $y = x + \frac{1}{2} \min (\varepsilon, 1)$. Then 
    \begin{itemize}
        \item $|y - x| = \frac{1}{2} \min (\varepsilon, 1) \implies y \in N(x, \varepsilon)$
        \item $x < y < x + 1 \implies y \notin \mathbb{N}$ 
    \end{itemize} 
    Hence $N(x, \varepsilon) \not \subset \mathbb{N}$ for all $\varepsilon > 0 \implies x \notin \operatorname{Int}  \mathbb{N}$.
\end{proof}

\begin{example}
    $\operatorname{bd} \mathbb{N} = \mathbb{N}$.
\end{example}
\begin{proof}
    \begin{enumerate}
        \item $\mathbb{N} \subset \operatorname{bd}  N$.
Fix $x \in \mathbb{N}, \varepsilon > 0$, 
        \begin{itemize}
            \item 
$x \in N(x, \varepsilon) \cap \mathbb{N} \neq \emptyset$. 
            \item 
$y = x + \frac{1}{2} \min(\varepsilon, 1) \in N(x, \varepsilon) \cap \mathbb{N}^c \neq \emptyset$.
        \end{itemize} 
        \item $\mathbb{N} \supset \operatorname{bd}  N$. Fix arbitrary $x \in \operatorname{bd} \mathbb{N}$. 
ATC $x \notin \mathbb{N}$, then we consider 2 cases
\begin{itemize}
    \item if $x < 1$, then take $\varepsilon = 1 - x$. 
$\forall y \in N(x, \varepsilon), 1 > y \notin \mathbb{N}$, 
hence $N(x, \varepsilon) \cap \mathbb{N} = \emptyset$, 
which contradicts the definition of boundary point. 
    \item $if x > 1$, let $n \in \mathbb{N}$ such that 
$n < x < n + 1$. Let $\varepsilon = \frac{1}{2}\min(x - n, n + 1 - x)$, then 
$\forall y \in N(x, \varepsilon)$, $y \notin \mathbb{N}$, 
hence $N(x, \varepsilon) \cap \mathbb{N} = \emptyset$, 
which contradicts the definition of boundary point. 
\end{itemize} 
    \end{enumerate}
\end{proof}

\begin{remark}
    \hfill 
    \begin{enumerate}
        \item $\operatorname{int} S = S \setminus \operatorname{bd} S$
        \item boundary points may or may not be in $S$. 
        \item Any point in a set $S$ is either a boundary or an interior point, but not both.
    \end{enumerate}
\end{remark}


















